\chapter{Testresultaten}
\label{app:testresultaten}

\section{Doel van de tests}
De tests evalueren of de proof of concept voldoet aan de functionele en
technische vereisten. De nadruk ligt op de werking van de editor, het
opslaan en laden van templates, de verwerking van HTML, het invullen van
placeholders en de weergave van gegenereerde e-mails.

Daarnaast wordt nagegaan of de oplossing ongewenste HTML-inhoud voldoende
afschermt en of de e-mailuitvoer bruikbaar is in gangbare e-mailclients.

\section{Testomgeving}
TODO!!!: Vul deze tabel aan met de werkelijk gebruikte configuratie.

\begin{table}[htbp]
    \centering
    \caption{Configuratie van de testomgeving}
    \label{tab:testomgeving}
    \begin{tabular}{|p{4cm}|p{9cm}|}
        \hline
        \textbf{Onderdeel} & \textbf{Configuratie} \\
        \hline
        Besturingssysteem & [In te vullen] \\
        \hline
        Frontend & Angular [versie invullen] \\
        \hline
        Backend & Node.js, TypeScript en Express [versies invullen] \\
        \hline
        Editor & [Geselecteerde editor en versie] \\
        \hline
        Testframework & [Bijvoorbeeld Vitest en Playwright, indien gebruikt] \\
        \hline
        E-mailtestomgeving & [Bijvoorbeeld Mailpit, indien gebruikt] \\
        \hline
        Testgegevens & Uitsluitend fictieve gegevens \\
        \hline
        Testdatum & [Datum invullen] \\
        \hline
    \end{tabular}
\end{table}

\section{Teststrategie}
De tests worden opgedeeld in verschillende categorieën. Niet elke categorie
is noodzakelijk al geïmplementeerd; de onderstaande indeling dient als
testplan en registratiesjabloon.

\begin{itemize}
    \item \textbf{Functionele tests:} controleren of de editor, preview,
    opslag en het laden van templates correct werken.
    \item \textbf{Renderingtests:} controleren of placeholders worden
    ingevuld en of de uiteindelijke HTML de verwachte structuur bevat.
    \item \textbf{Veiligheidstests:} controleren of niet-toegestane HTML,
    attributen en scripts worden verwijderd of geweigerd.
    \item \textbf{Geautomatiseerde tests:} controleren de werking van
    afzonderlijke functies en belangrijke gebruikersflows.
    \item \textbf{Compatibiliteitstests:} controleren hoe de gegenereerde
    e-mails worden weergegeven in beschikbare e-mailclients.
\end{itemize}

\section{Functionele testgevallen}

TODO!!!:Registreer elk testgeval met de werkelijk verkregen uitkomst.

\begin{table}[htbp]
    \centering
    \caption{Functionele testgevallen}
    \label{tab:functionele-tests}
    \small
    \begin{tabular}{|p{1.2cm}|p{3cm}|p{4.2cm}|p{3cm}|p{1.5cm}|}
        \hline
        \textbf{ID} & \textbf{Test} & \textbf{Verwachte uitkomst} &
        \textbf{Werkelijke uitkomst} & \textbf{Status} \\
        \hline
        F-01 & Template laden &
        De geselecteerde template wordt correct weergegeven. &
        [Invullen] & [Open/Geslaagd/Mislukt] \\
        \hline
        F-02 & Tekst bewerken &
        De gebruiker kan de inhoud wijzigen. &
        [Invullen] & [Invullen] \\
        \hline
        F-03 & Opmaak aanpassen &
        Ondersteunde tekstopmaak wordt toegepast. &
        [Invullen] & [Invullen] \\
        \hline
        F-04 & Template opslaan &
        De aangepaste template kan worden opgeslagen. &
        [Invullen] & [Invullen] \\
        \hline
        F-05 & Template opnieuw laden &
        De opgeslagen wijzigingen blijven behouden. &
        [Invullen] & [Invullen] \\
        \hline
        F-06 & Preview openen &
        De preview toont de gerenderde e-mail. &
        [Invullen] & [Invullen] \\
        \hline
        F-07 & Testmail versturen &
        De e-mail verschijnt in de lokale testomgeving. &
        [Invullen] & [Invullen] \\
        \hline
    \end{tabular}
\end{table}

\section{Rendering- en veiligheidstests}

\begin{table}[htbp]
    \centering
    \caption{Testgevallen voor rendering en veiligheid}
    \label{tab:veiligheidstests}
    \small
    \begin{tabular}{|p{1.2cm}|p{3cm}|p{4.2cm}|p{3cm}|p{1.5cm}|}
        \hline
        \textbf{ID} & \textbf{Test} & \textbf{Verwachte uitkomst} &
        \textbf{Werkelijke uitkomst} & \textbf{Status} \\
        \hline
        R-01 & Geldige placeholder &
        Een toegestane placeholder wordt correct ingevuld. &
        [Invullen] & [Invullen] \\
        \hline
        R-02 & Onbekende placeholder &
        Een onbekende placeholder wordt veilig afgehandeld. &
        [Invullen] & [Invullen] \\
        \hline
        R-03 & HTML-sanitatie &
        Niet-toegestane HTML wordt verwijderd of geweigerd. &
        [Invullen] & [Invullen] \\
        \hline
        R-04 & Schadelijke attributen &
        Niet-toegestane attributen, zoals event-handlers, worden niet uitgevoerd. &
        [Invullen] & [Invullen] \\
        \hline
        R-05 & Ontbrekende gegevens &
        Ontbrekende gegevens veroorzaken geen onveilige of misleidende uitvoer. &
        [Invullen] & [Invullen] \\
        \hline
        R-06 & CSS-inlining &
        Ondersteunde stijlen worden in de HTML-uitvoer verwerkt. &
        [Invullen] & [Invullen] \\
        \hline
        R-07 & Tekstversie &
        Indien voorzien, wordt een bruikbare tekstversie gegenereerd. &
        [Invullen] & [Invullen] \\
        \hline
    \end{tabular}
\end{table}

\section{Compatibiliteit met e-mailclients}
E-mailclients kunnen HTML en CSS verschillend weergeven. De compatibiliteit
wordt daarom gecontroleerd aan de hand van dezelfde testtemplates en
fictieve gegevens.

\begin{table}[htbp]
    \centering
    \caption{Compatibiliteitstest per e-mailclient}
    \label{tab:compatibiliteit}
    \small
    \begin{tabular}{|p{3cm}|p{2.5cm}|p{2.5cm}|p{3.5cm}|p{2cm}|}
        \hline
        \textbf{E-mailclient} & \textbf{Testdatum} &
        \textbf{Uitvoering} & \textbf{Vastgestelde afwijkingen} &
        \textbf{Status} \\
        \hline
        Gmail & [Datum] & [Versie/platform] & [Invullen] & [Invullen] \\
        \hline
        Outlook.com / Outlook & [Datum] & [Versie/platform] &
        [Invullen] & [Invullen] \\
        \hline
        Apple Mail & [Datum] & [Mac/iPhone, versie] &
        [Invullen] & [Invullen] \\
        \hline
    \end{tabular}
\end{table}

TODO!!!:Voer de tests uit in de werkelijk beschikbare clients. Indien een client niet
beschikbaar is, vermeld dan dat de compatibiliteit voor die client niet werd
geverifieerd.

\section{Automatische testresultaten}

TODO!!!:Vul deze tabel in met de uitvoer van de gebruikte testtools.

\begin{table}[htbp]
    \centering
    \caption{Samenvatting van de automatische tests}
    \label{tab:automatische-tests}
    \begin{tabular}{|p{4cm}|p{2cm}|p{2cm}|p{2cm}|p{2.5cm}|}
        \hline
        \textbf{Testcategorie} & \textbf{Uitgevoerd} &
        \textbf{Geslaagd} & \textbf{Mislukt} & \textbf{Opmerkingen} \\
        \hline
        Unit tests & [Aantal] & [Aantal] & [Aantal] & [Invullen] \\
        \hline
        Integratietests & [Aantal] & [Aantal] & [Aantal] & [Invullen] \\
        \hline
        Browsertests & [Aantal] & [Aantal] & [Aantal] & [Invullen] \\
        \hline
    \end{tabular}
\end{table}

\section{Vastgestelde problemen en oplossingen}

TODO!!!:Beschrijf hier de belangrijkste problemen die tijdens het testen werden
vastgesteld. Noteer per probleem de oorzaak, de genomen maatregel en het
resultaat van een eventuele hertest.

\begin{table}[htbp]
    \centering
    \caption{Overzicht van testproblemen}
    \label{tab:testproblemen}
    \begin{tabular}{|p{1.5cm}|p{3cm}|p{3cm}|p{3cm}|p{2cm}|}
        \hline
        \textbf{ID} & \textbf{Probleem} & \textbf{Oorzaak} &
        \text{Oplossing} & \textbf{Hertest} \\
        \hline
        P-01 & [Invullen] & [Invullen] & [Invullen] & [Invullen] \\
        \hline
        P-02 & [Invullen] & [Invullen] & [Invullen] & [Invullen] \\
        \hline
    \end{tabular}
\end{table}

\section{Conclusie}

TODO!!!:Op basis van de uitgevoerde tests wordt beoordeeld in welke mate de PoC
voldoet aan de vooropgestelde vereisten.

De belangrijkste positieve resultaten zijn: [invullen].
De resterende beperkingen zijn: [invullen].
De volgende vereisten zijn volledig aangetoond: [invullen].
De volgende vereisten konden niet of slechts gedeeltelijk worden aangetoond:
[invullen].

Deze conclusies zijn gebaseerd op de geregistreerde testresultaten en gelden
voor de onderzochte PoC en testomgeving. Ze vormen geen garantie voor de
werking van een toekomstige productie-integratie.